\documentclass[10pt,a4paper]{amsart}
\usepackage[margin=19mm]{geometry}
\usepackage{amsmath,amssymb,mathtools}
\usepackage[scr=rsfs,cal=euler]{mathalfa}
\usepackage{xfrac}
\usepackage[unicode,hidelinks,bookmarks=false]{hyperref}
\newcommand{\ie}{i.e.\ }
\newcommand{\cf}{cf.\ }
\newcommand{\resp}{respectively\ }
\newcommand{\Subsection}{\subsection}
\providecommand{\nobreakdash}{\nobreak\hbox{-}}
\setlength{\emergencystretch}{2em}
\multlinegap0pt
\usepackage{amssymb}
\usepackage{mathtools}
\usepackage{tikz}
\usepackage{multirow}
\newtheorem{theorem}{Theorem}
\title{A statistic-swapping involution on the Cartesian product of the symmetric group \texorpdfstring{$S_{kn}$}{on kn letters} and the generalized symmetric group \texorpdfstring{$S(k,n)$}{S(k,n)}}
\author{Peter Kagey, Kai Mawhinney}
\date{Source: arXiv:2602.11061v1 (CC BY 4.0)}
\begin{document}
\maketitle

\noindent\emph{Source and licence.} A statistic-swapping involution on the Cartesian product of the symmetric group \texorpdfstring{$S_{kn}$}{on kn letters} and the generalized symmetric group \texorpdfstring{$S(k,n)$}{S(k,n)}. Version arXiv:2602.11061v1 (CC BY 4.0), distributed by arXiv under the Creative Commons Attribution 4.0 International licence, http://creativecommons.org/licenses/by/4.0/, as recorded in the arXiv licence metadata for this exact version and shown on the abstract page https://arxiv.org/abs/2602.11061v1. Full licence notice: \texttt{licenses/arxiv-2602.11061-CC-BY-4.0.txt}.\par\smallskip

\noindent\textit{Editorial scope.} Counted unit: the article's theorem thm:involution (source0.tex lines 334--342). The article's complete original proof of it is delivered with it (source0.tex lines 343--354, one \texttt{proof} environment of 12 lines). No mathematics is written, repaired or re-derived: the statement and the proof are the article's own lines, in the article's own notation, and no retained line is substituted or re-flowed.  No further source-local context is needed: every cross-reference of every retained line resolves inside the two delivered spans.  Nothing was omitted from any retained span, so the package carries no omission record.  No retained line contains a citation, so the wrapper carries no bibliography and no bibliography entry.  No retained line contains a figure environment, an image-inclusion command or drawing code, so no image is delivered and the package declares no asset dependency.  Editorial formatting only, in addition: an AMS \texttt{amsart} preamble that restates the article's own declarations and macros used by the retained lines, namely the environments \texttt{theorem} on separate counters, together with the standard packages those lines need.  The retained environments print in this document as Theorem 1 in order of appearance, whereas the article prints the same items with the numbering of its own full document.  The complete native source of the article as distributed in the arXiv e-print for this exact version is bundled unchanged as \texttt{sources/194459/source0.tex}.
\begin{theorem}\label{thm:involution}
    Suppose that $f\colon S_{kn} \to D_{k,n} \times S(k,n)$ is a bijection such that $f(\pi) = \left(f_1(\pi), f_2(\pi)\right) = (\delta, \sigma)$ and the number of $k$-cycles in $\pi$ is equal to the number of fixed points in $\sigma$.

    Then the map $\phi \colon S(k,n) \times S_{kn} \to S(k,n) \times S_{kn}$ defined as \[
        \phi(\sigma', \pi) = \left(f_2(\pi),f^{-1}(f_1(\pi),\sigma')\right)
        = \left(\sigma, f^{-1}(\delta, \sigma')\right)
    \]
    is an involution, such that the number of $k$-cycles in $\pi$ is equal to the number of fixed points in $\sigma$ and the number of $k$-cycles in $\pi'$ is equal to the number of fixed points in $\sigma'$.
\end{theorem}

\begin{proof}
    Let $(\sigma', \pi) \in S(k,n) \times S_{kn}$. Denote $f(\pi) = (\delta, \sigma)$ and denote $f^{-1}(\delta, \sigma') = \pi'$.
    To check that this is an involution, we apply the map twice: \begin{align*}
        \phi(\phi(\sigma', \pi))
        &= \phi\!\left(\sigma, f^{-1}(\delta, \sigma')\right)
        = \phi(\sigma,\pi') \\
        &= \left(f_2(\pi'),f^{-1}(f_1(\pi'),\sigma)\right)
        = \left(\sigma',f^{-1}(\delta,\sigma)\right)
        = \left(\sigma', \pi\right)
    \end{align*}
    Moreover, by the assumption on $f$, if $\pi$ has $\alpha$ $k$-cycles, then $\sigma'$ has $\alpha$ fixed points, and if $\pi'$ has $\beta$ $k$-cycles, then $\sigma$ has $\beta$ fixed points. Thus $\phi(\sigma, \pi) = (\sigma', \pi')$ switches the number of $k$-cycles and fixed points.
\end{proof}

\end{document}
